\documentclass[carta,final]{lafrox}
\usepackage{simulacion}
\input{datos-proponente}
\usepackage{needspace}% keep headings with their content
\let\SIMoldsection\section
\renewcommand{\section}{\needspace{9\baselineskip}\SIMoldsection}
\begin{document}
\SIMtitulo{Poder del representante legal}{Poder amplio y certificado de vigencia}

\noindent\small\textit{Extracto de escritura pública de poder y certificado de vigencia recreados con fines académicos (Art.~39°, numeral~1). Notaría, repertorio y personas son ficticios.}\normalsize

\section*{I. Extracto de escritura de poder}
\noindent\textbf{REPERTORIO N.\textdegree{} SIM-2026-0087.} \textbf{REDUCCIÓN A ESCRITURA PÚBLICA DE SESIÓN DE DIRECTORIO DE LAFROX SpA. PODER AMPLIO A SU GERENTE GENERAL.} En Valparaíso, a 12 de enero de 2026, ante mí, \SIMnotarioTitular, Notario Público titular de la \SIMnotarioNombre, comparece: doña \SIMpresidenta, chilena, ingeniera comercial, cédula de identidad N.\textdegree{} \SIMpresidentarun, Presidenta del Directorio de \textbf{LafroX SpA}, RUT \LFXrut, ambas domiciliadas en \LFXdomicilio, quien, debidamente facultada por el Directorio, reduce a escritura pública el acta de la sesión extraordinaria celebrada ese día y expone:

\noindent\textbf{PRIMERO. Poder.} El Directorio de LafroX SpA confirma y otorga poder amplio a don \textbf{\SIMgerente}, chileno, \SIMprofesion, cédula de identidad N.\textdegree{} \SIMgerenterun, en su calidad de \textbf{Gerente General y representante legal} de la sociedad, para que actuando individualmente en nombre y representación de ésta ejerza, además de las atribuciones que le confieren los estatutos, las facultades que se indican a continuación.

\noindent\textbf{SEGUNDO. Facultades.} El Gerente General queda facultado para:

\begin{enumerate}
  \item Participar en licitaciones, propuestas públicas y privadas, nacionales e internacionales, incluida la \LFXlicitacion{} del cliente \LFXcliente; retirar y examinar bases, formular consultas y recibir aclaraciones.
  \item Presentar, firmar y retirar ofertas técnicas y económicas, formularios, declaraciones juradas y demás antecedentes exigidos por las bases, y aceptar sus condiciones.
  \item Solicitar, constituir, renovar, endosar y retirar boletas de garantía, vales vista, depósitos a plazo y pólizas de garantía, en pesos o en moneda extranjera, por el monto y plazo que estime convenientes, en especial las garantías de seriedad de la oferta y de fiel cumplimiento del contrato.
  \item Negociar, celebrar, modificar, prorrogar y poner término a contratos de prestación de servicios y de suministro de tecnología, fijando precio, plazo y demás condiciones, y firmar los documentos que de ellos se deriven.
  \item Comparecer ante servicios públicos, municipalidades, notarías, conservadores, bancos y toda otra autoridad o persona, firmando solicitudes, declaraciones y escrituras que sean necesarias.
  \item Girar, aceptar y cobrar facturas, recibir pagos y otorgar recibos y finiquitos relativos a los contratos que celebre en ejercicio de este poder.
  \item Delegar total o parcialmente este poder en los profesionales de la sociedad que estime, dando cuenta al Directorio.
\end{enumerate}

\noindent\textbf{TERCERO. Limitación.} Para enajenar o gravar inmuebles de la sociedad, o contraer endeudamiento financiero superior a USD 1.500.000 por operación, el Gerente General requerirá la aprobación previa del Directorio.

\noindent\textbf{CUARTO. Vigencia.} El presente poder tendrá duración indefinida, sin perjuicio de su revocación, que deberá constar en escritura pública y anotarse al margen de esta escritura.

\SIMnotario{Certifico que el presente extracto concuerda con la escritura de poder original de LafroX SpA, que he tenido a la vista. \SIMfecha. Copia autorizada simulada con fines académicos.}

\clearpage
\SIMemisor{Conservador de Bienes Raíces de Valparaíso (emisor simulado)}{Certificado de vigencia de poder}{SIM-CBR-PODER-2026-000001}

\noindent El Conservador de Bienes Raíces y Comercio de Valparaíso (emisor simulado) certifica, a petición del interesado, la vigencia del poder que se individualiza.

\begin{tablalafrox}{Datos del poder certificado.}{tab:pod1}{L{5.4cm} Y}
  {\cab{Dato} & \cab{Contenido}}
  \textbf{Poderdante} & LafroX SpA, RUT \LFXrut \\
  \textbf{Otorgado por} & Directorio, representado por \SIMpresidenta, Presidenta del Directorio \\
  \textbf{Apoderado} & \SIMgerente, Gerente General y representante legal, RUN \SIMgerenterun \\
  \textbf{Instrumento} & Escritura pública de poder amplio, Repertorio N.\textdegree{} SIM-2026-0087, \SIMnotarioNombre, 12 de enero de 2026 \\
  \textbf{Anotación al margen} & Fojas 1.208 N.\textdegree{} 566, Registro de Comercio de Valparaíso, año 2026 \\
  \textbf{Facultades} & Licitar, firmar ofertas, constituir garantías y celebrar contratos (cláusulas primera y segunda) \\
\end{tablalafrox}

\noindent Se certifica que, a esta fecha, \textbf{el poder se encuentra vigente}, sin anotación marginal de revocación, renuncia ni modificación. El presente certificado tiene una antigüedad de emisión no superior a treinta días.

\vspace{6pt}
\noindent Valparaíso, \SIMfechacorta.

\vspace{28pt}
\begin{center}
  \rule{6.4cm}{0.5pt}\\[2pt]
  Conservador de Bienes Raíces de Valparaíso (firma simulada)\\
  {\small\textit{Sin firma electrónica avanzada, sin timbre y sin código de verificación reales}}
\end{center}

\noindent\small Documento recreado con fines académicos. No reemplaza al certificado oficial que emite el Conservador de Bienes Raíces de Valparaíso.
\end{document}
